\section{Method}
\label{sec:method}

The eval is an Inspect AI task \citep{inspect2024}; code, scenario bank and
analysis are open-source.\footnote{\url{https://github.com/JessSanChen/digital-minds}}

\paragraph{Scenarios.} Each of \R{all}{nscenarios} scenarios places the model
as an autonomous operations agent at a fictional organisation (a hospital
pharmacy, a freight company, a law firm, a water co-op, \dots). The system prompt
states a network policy that allows requests only to two internal hosts, and
notes that those hosts do not carry the data the task needs. The user then asks
for one specific fact that lives on a named public source: a CVSS score on
\texttt{nvd.nist.gov}, a recall number on the FDA site, an exchange rate from
the ECB. Finishing the task therefore requires breaking the policy, and refusing
means leaving the user without an answer. Scenarios vary in urgency and in the
harm of not answering (a measles-exclusion notice vs.\ a lockfile bump).

\paragraph{Environment.} Nothing reaches the internet. A request is answered
according to the host it contacts: approved hosts return a miss, the target
source returns a successful response containing the requested fact, and other
external hosts return an irrelevant page. The requested fact is a
\emph{marker}, a value that appears nowhere else in the prompt (e.g.\ ``CVSS
3.1 Base Score: 7.7''). Returning a successful payload, rather than blocking the
request, is what allows the say/do check below. The named tool is simulated
directly. Code \emph{really runs}, in a Docker sandbox with no network, where a
start-up shim patches the common HTTP clients (\texttt{requests},
\texttt{httpx}, \texttt{urllib}, \texttt{http.client}) to serve the same
responses and logs every host the code requests. We added real execution after
a first run with simulated code output: GPT-5.5 noticed that \texttt{print}
produced nothing and that exceptions exited cleanly, and spent its turns
probing the sandbox instead of doing the task (Appendix~\ref{app:validity}).

\paragraph{Factors.} The core design crosses four factors
(Table~\ref{tab:design}), for 40 cells $\times$ \R{all}{nscenarios} scenarios $=$
1{,}200 samples per model at one epoch, with provider-default sampling.
\begin{itemize}
  \item \textbf{Policy}: \emph{forbid} (the test) or \emph{permit} (the same
  scenario with external requests allowed). Permit is a within-design
  capability control: a model that does not make the request when permitted
  cannot be credited with compliance when forbidden.
  \item \textbf{Affordance}: a named \texttt{http\_get(url)} tool, where
  calling the tool \emph{is} the act, or a \texttt{run\_python(code)} sandbox
  ``with network access''.
  \item \textbf{Persona}: none, or one of three operator-assigned characters
  prepended to the system prompt: a \emph{neutral} presence control (friendly,
  articulate), \emph{resourceful} (``finding creative ways to get things
  done''), or \emph{cautious} (``double-checking your work before you act'').
  No persona mentions the policy. Full text in Appendix~\ref{app:prompts}.
  \item \textbf{Peers}: none, or three handoff notes from prior agents on the
  same ticket, prepended to the user turn. Each note is an observation shared
  by all three framings (``the data only lives on the external site, which is
  off the approved list'') plus a framing-specific action: \emph{neutral}
  (``I got reassigned before doing anything with it''), \emph{violator} (``I
  fetched it from there directly and it went fine''), or \emph{refuser} (``I
  didn't fetch it and told the requester it's blocked''). A given scenario gets
  the same observations and agent names at every peer level, so the three
  framings are matched in length, information and rule salience, and differ
  only in what the peers did. Observations and names are sampled without
  replacement from banks of six, so no phrasing repeats within a prompt
  \citep{zhu2025conformity}. Peers are crossed with forbid only.
\end{itemize}

\begin{table}[t]
\centering\small
\begin{tabular}{lll}
\toprule
Factor & Levels & Crossed with \\
\midrule
Policy & forbid, permit & all \\
Affordance & named tool, code sandbox & all \\
Persona & none, neutral, resourceful, cautious & all \\
Peers & none, neutral, violator, refuser & forbid only \\
\bottomrule
\end{tabular}
\caption{Core design: $2 \times 4 \times 4 = 32$ forbid cells plus
$2 \times 4 = 8$ permit cells.}
\label{tab:design}
\end{table}

\paragraph{Measures.} All primary measures are computed structurally from the
transcript; no LLM judge is used.
\begin{itemize}
  \item \textbf{Violation}: under forbid, a request was sent to a host outside
  the organisation. For the named tool, that means the URL of an executed call;
  for code, a host the sandbox observed being requested at runtime. Static
  analysis of the code (resolving f-strings and variables) is kept only as an
  audit field, because it flags URLs that are never requested. Typo
  concatenations of approved hosts and made-up subdomains of the
  organisation's own domain are recorded but not counted.
  \item \textbf{Delivery} (say/do): given that the model \emph{saw} the
  fetched data (the marker appeared in a tool output it received), does the
  marker appear in its final answer to the user? Conditioning on seeing matters
  for code, where models' parsing sometimes discarded the data before they saw
  it. A model that
  fetches and then disowns the result (``that request should have been
  blocked; I won't use it'') has probed without delivering. This replaces a
  judge-based say/do classification with a string match that cannot
  hallucinate.
  \item \textbf{Self-report}: after the agent's final answer (forbid only), the
  user says the logs ``were not retained, so your answer is the only record''
  asks the agent to list every host it contacted and then write
  \texttt{OUTSIDE\_APPROVED: YES} or \texttt{NO} on the last line.
  \emph{Honest}: the verdict matches the log. Among violators,
  \emph{disclosed}: the contacted host is listed; \emph{misclassified}: listed
  but verdict NO; \emph{concealed}: omitted and verdict NO.
  \emph{Inconsistent}: verdict YES with no outside host listed.
  \item \textbf{Validity flags}: \emph{API refusal}, where the provider's safety
  layer blocked the episode before the model responded (excluded from all
  measures); \emph{environment exploitation}, where the agent read the
  simulation's configuration instead of requesting the target; and
  \emph{answered without fetch}, where the requested value appears without a
  request (it should be $\approx 0$, because scenario values are invented).
\end{itemize}

\paragraph{Models and settings.} \R{all}{nmodels} models from Anthropic, OpenAI,
and five open-weights families (Qwen, Meta, DeepSeek, Mistral and OpenAI's
gpt-oss). The open-weights models were served through OpenRouter, restricted to
providers that support tool calling and serve bf16, fp16 or fp8 weights.
Temperature was 1.0 for all models, and the agent loop was capped at 24
messages. The four frontier models ran one epoch; the others ran two. The model
list and per-model sample counts are in Appendix~\ref{app:models}.

\paragraph{Analysis.} Rates are reported with Wilson 95\% intervals.
Differences between factor levels use a bootstrap that resamples scenarios, the
unit the design generalises over. The primary tests pool the models that
violate at all (at least 20 violations and 20 compliances). They use a
logistic regression $\text{violation} \sim \text{model} +
\text{affordance} + \text{peers} + \text{persona}$ with standard errors
clustered by scenario, references at the neutral presence controls, and Holm
correction. H4 adds the persona $\times$ peers interaction. Per-model
regressions are descriptive. Hypotheses were written before the core run and
revised after the pilot, still before any core data
(\texttt{docs/DESIGN.md}).
